% Terminal note of the homily: everything the speech itself must not carry ---
% the audience, the spoken word count and the pace it implies, the relation to
% the three interpretations of the expansive study, the connections that are
% the homily's own, the route by which each quoted text reaches the page, and
% the exact loci --- followed by the References. Nothing here is read aloud.

\section*{Note on Sources and Delivery}
\runninghead{Sources and Delivery}

{\footnotesize

\textbf{Audience and occasion.} An adult parish assembly, at the Mass
\latin{In festo omnium Sanctorum}, I classis, of the 1962 \work{Missale
Romanum}, editio typica, pp.~702--703, nos.~3974--3983, said on Sunday
1~November 2026 with the commemoration of the Twenty-third Sunday after
Pentecost. The speech presumes that it follows the Gospel, and it speaks of
the Offertory, Secret, Communion and Postcommunion as still to come. It makes
no use of the Sunday's three commemorated orations, the Alleluia or the
Preface.

\textbf{Length and pace.} The spoken body runs to about 1,370 words, which is
about 10.5 to 11.4 minutes at an unhurried 120 to 130 words a minute, by
arithmetic alone; no one has delivered or timed it. It was read through in
full, silently, for sense and ease of speech.

\textbf{Relation to the reviewed interpretations.} The homily follows chiefly
the third of the expansive study's three interpretations, \textit{No Want to
Them That Fear Him, and the Just in God's Hand}: Augustine's hearer who
complains that the Gradual's promise has failed him, and his answer that the
goods are spiritual; Basil's dilemma from the apostles, his ``God himself''
as the good the seeker cannot lack, and his fear that lacks nothing for any
virtue; Aquinas's help in need and the physician, with the difference between
Basil and Aquinas over bodily goods kept; Augustine's gloss on the Offertory's
first clause; and Bernard on the verse as sung in the night office. From the
second, \textit{The Beatitudes, the Way the Saints Walked}, it takes the
present tense of the first and eighth beatitudes, Augustine's poor in spirit
as the humble and his pairing of them with the fear of the Lord, and
Chrysostom's ``Both here and there''; Hilary's emphasis on the life after
death, Bernard's voluntary poverty as a kind of martyrdom, and the others'
authorities on the tense are left to the studies. From the first, \textit{One
City of Angels and Men, Numbered by God}, it takes only the Lesson's
multitude, the Collect's and Postcommunion's prayer through the saints, and
Bernard on the veneration of the saints. Bellarmine, the dispute over whose
souls the Offertory names, and the number of the saved are left to the
studies; the speech says ``the just,'' as the antiphon does, and does not
narrow it to the martyrs. Two details go beyond what the studies report, and
each was read at its locus. Augustine judges that whoever reasons as his
hearer does is caught in the snare of scandal, seeking perishable food on
earth and not the true reward in heaven, and then tells him not to
understand the verse so (Migne Latin, nn.~14--15). And the physician with
which Aquinas meets the objection from St~Paul's hunger is, on Aquinas's own
report, Augustine's: Aquinas says that Augustine raises the objection and
that God is nourisher and physician, and Augustine, on Mt 6:33, answers
St~Paul's hunger with the Physician who gives and takes away such things as he
judges expedient for us.

\textbf{Connections that are the homily's own.} These are drawn in the
preacher's voice and credited to no Father. The Gradual and the Gospel are
joined through Augustine's pairing: Augustine gives the fear of the Lord to
the humble, the poor in spirit, and the homily infers that the Gradual's man
who fears the Lord and the Gospel's poor in spirit are the same man, to whom
the Gradual promises that nothing is lacking and the first beatitude's
present says that the kingdom is his now; that the kingdom belongs to those
persecuted for justice' sake while they suffer it is read from the present
tense of Mt 5:10. The Gradual and the Offertory are taken together, so that
the just man put to death is said to have kept the good the Gradual
promises, Basil's ``God himself'' (the third interpretation's own
inference). The Secret's \latin{iustorum} is taken for the saints of the
feast. The Lesson's multitude is said to include the Gospel's poor in spirit,
mourners and persecuted and the Offertory's just. The Gradual's ``all ye his
saints'' is said to be spoken to the assembly, from Bernard's words on the
night office, which do not comment on the Mass Gradual; Basil's account of
the promise by the fear that keeps from every shameful act, which is his own
comment on Ps 33:10, follows it. The Collect and Postcommunion are said to
ask for us because, on Bernard's account, the saints' veneration serves us.
The response takes the daily practice of Augustine's hearer as its pattern,
and his measuring of what he had against the rich man's goods (Augustine's
\latin{ego timeo Dominum, ille non timendo quantum acquisivit}) as what it
corrects; Augustine does not offer the practice as a pattern, and the speech
reports his verdict on the man before the response is given. The Lord who is
sought for himself is said to give himself at the altar in Christ under the
sacrament, and the Communion's ``they shall see God'' is read of the vision
face to face (1~Cor 13:12, in the words of the second interpretation's
witnesses on Mt 5:8); Basil does not say this of his ``God himself.''
Basil's comment on Ps 32:1, the Introit's verse, is read with the Introit
antiphon's opening words.

\textbf{The half-verse the Gradual omits.} The speech quotes Ps 33:11a, ``The
rich have wanted, and have suffered hunger,'' from the Douay, and says that
the Gradual leaves it out. Augustine's hearer complains on the whole of
verse~11, the half the Gradual sings with the half it omits, and Augustine's
admission that the verse seems to deceive is said of the whole verse. His
comment on the omitted half, that the rich lack the living bread, is not
used.

\textbf{Texts quoted, and their route.} Scripture is the Douay--Rheims,
Challoner's revision: Ps 33:10 and 11a; Wis 2:20; Mt 5:3 and 10, and 5:8 at
the Communion; Apoc 7:9--10. The Gradual's second line (Ps 33:11b,
\latin{non deficient omni bono}), the futures of Mt 5:5, 6 and 8 in the
Gospel paragraph, and St~Paul's hunger in 1~Cor 4:11 are given in the
speech's own words. The Introit's first words, the Collect's ``through the
multitude of our intercessors,'' the Secret's ``for the honor of all the
just'' and the Postcommunion's ``defended by their unceasing prayers'' are
the English of the Lasance \work{New Roman Missal} (1945), a historical study
aid and not an approved liturgical translation of the 1962 Missal; the
Offertory is the Lasance's English of the antiphon, quoted without its
closing \latin{Alleluia}. Three patristic passages are quoted, each in
quotation marks in the speech and each in the English the expansive study
gives beside its Latin: Augustine's hearer's complaint (\work{Enarr.\ in
Ps.}\ 33), Augustine's gloss on the just in God's hand, ``dwell in the help
of the Most High, and abide in the protection of the God of heaven''
(\work{Enarr.\ in Ps.}\ 87), and Bernard's ``Meanwhile our whole beatitude
is to fear God.'' Chrysostom's ``Both here and there'' is quoted in the
English of the Nicene and Post-Nicene Fathers. Everything else from the
Fathers is paraphrased: Augustine's admission and the snare, from the Migne
Latin of n.~14, and his telling the hearer not to understand the verse so and
the spiritual goods, from n.~15; Augustine on the Sermon on the Mount,
including the physician of II.17.58, from the Nicene and Post-Nicene English;
Basil throughout; Bernard on the saints in glory who no longer fear, on their
feast in joy and praise, and on the veneration that serves us; and Aquinas.

\textbf{Exact loci.} Introit, no.~3974; Collect, no.~3975; Lesson, Apoc
7:2--12, no.~3976; Gradual, Ps 33:10, 11b (34:10, 11b), no.~3977; Gospel,
Mt 5:1--12a, no.~3979; Offertory, Wis 3:1--3, adapted, no.~3980; Secret,
no.~3981; Communion, Mt 5:8--10, no.~3982; Postcommunion, no.~3983.
Augustine, \work{Enarr.\ in Ps.}\ 33, s.~2, n.~14 (the hearer's complaint;
\latin{ego timeo Dominum, ille non timendo quantum acquisivit}; \latin{et
dicis in animo tuo}; the verse seeming to deceive; \latin{Sic cogitantem
suffocat laqueus scandali. Escam enim mortalem quaerit in terra, et veram
mercedem non quaerit in coelo}) and n.~15 (\latin{Noli ergo sic intelligere.
Et quomodo intelligam? In bonis spiritualibus}), and 87, n.~5 (the just in
God's hand, with Ps 90:1); \work{De serm.\ Dom.}\ I.3.10 and I.4.11 (the poor
in spirit as the humble, not puffed up, who submit to God; the gifts counted
upward, ``in a gradually ascending series,'' so that the fear of God is
first and ``corresponds to the humble''), and II.17.58 (St~Paul's hunger and
thirst; ``that Physician'' who ``knows such things just as helps,'' ``when He
takes them away, just as He judges it expedient for us''). Basil, \work{Hom.\ in
Ps.}\ 33, n.~6 (\latin{Timenti itaque non est inopia}: the fearer kept from
every unseemly act lacks nothing for any virtue; the speech says shameful)
and n.~7 (the dilemma; St~Paul; the perfect good), and \work{Hom.\ in Ps.}\
32, n.~1 (rejoicing because they have the Lord). Aquinas, \work{In Ps.}\ 33,
n.~10 (the spiritual goods; help in need; \latin{Sed objicit Augustinus in
serm.\ Dom.\ in mont.}, with 1~Cor 4:11; \latin{Et dicit quod Deus est
nutritor et medicus}). Chrysostom, \work{Hom.\ in Matt.}\ 15.4 (the mourners
comforted).
Bernard, \work{In festo omnium Sanctorum}, sermo~5, n.~3 (the verse sung in
the night office; ``our whole beatitude''; \latin{nostra interim festivitas,
in timore Dei: nam illorum magis in exsultatione et laude versatur}) and n.~5
(the saints do not need our honours; veneration concerns us, not them, and
kindles the desire for their company).

\par}

\Needspace{14\baselineskip}
\section*{References}
\runninghead{References}

{\footnotesize
\begin{itemize}\setlength{\itemsep}{0.15em}\setlength{\parskip}{0pt}
\item \work{Missale Romanum ex decreto Sacrosancti Concilii Tridentini restitutum, Summorum Pontificum cura recognitum}, editio typica (Typis Polyglottis Vaticanis, 1962), CMAA facsimile, Sanctorale, pp.~702--703, nos.~3974--3983.
\item Lasance, F.~X., \work{The New Roman Missal in Latin and English}, revised edition (New York: Benziger Brothers, 1945), pp.~1265--1268 (All Saints: Introit, Collect, Offertory, Secret, Postcommunion).
\item \work{The Holy Bible Translated from the Latin Vulgate} (Douay--Rheims, Challoner revision), Project Gutenberg text: Ps 33:10, 11a; Wis 2:20; Mt 5:3, 8, 10; Apoc 7:9--10.
\item St Augustine, \work{Enarrationes in Psalmos} 33 (sermo~2, nn.~14--15) and 87 (n.~5), Migne Latin text (Wikisource); \work{De sermone Domini in monte} I.3.10, I.4.11 and II.17.58, Nicene and Post-Nicene Fathers, 1st ser., vol.~6.
\item St Basil the Great, \work{Homilia in Psalmum XXXII}, n.~1, and \work{Homilia in Psalmum XXXIII}, nn.~6--7, \work{Patrologia Graeca} 29 (Paris, 1860), Latin version.
\item St Bernard of Clairvaux, \work{In festo omnium Sanctorum} sermo~5, nn.~3 and 5, \work{Patrologia Latina} 183 (Paris, 1854).
\item St John Chrysostom, \work{Homilies on Matthew} 15, Nicene and Post-Nicene Fathers, 1st ser., vol.~10.
\item St Thomas Aquinas, \work{In Psalmos Davidis expositio}, on Ps 33, in \work{Opera omnia}, vol.~14 (Parma, 1863), pp.~206--207; paraphrased, not quoted.
\end{itemize}
\par}
